% !TEX root = main.tex
\chapter{Uniqueness of Ramsey Expansions with the Ordering Property}

In Chapter~8 we saw that if a reasonable Fra\"iss\'e order expansion $\k$ of a Fra\"iss\'e class $\k_0$ has both the Ramsey property and the ordering property, then the flow $X_\k$ is the universal minimal flow of $\Aut(\mathrm{Flim}(\k_0))$. Since the universal minimal flow is unique up to isomorphism, one expects that such expansions should themselves be unique in an appropriate sense. The aim of this chapter is to make that statement precise.

We first introduce the notion of simple bi-definability between ordered expansions and prove the main uniqueness theorem. We then discuss some strengthenings in special cases and conclude with a minimality result showing that, up to simple bi-definability, expansions with both the Ramsey and ordering properties are the smallest Ramsey expansions of a given reduct class.

\section{Simple bi-definability and the main uniqueness theorem}

Let $L_0$ be a relational signature, and let
\[
    L=L_0\cup\{<\},
    \qquad
    L'=L_0\cup\{<'\},
\]
where $<$ and $<'$ are new binary relation symbols. In this section we compare two Fra\"iss\'e order expansions of the same reduct class
\[
    \k_0.
\]

We begin by formalizing the notion that one of the two orderings can be recovered from the other by a formula of particularly simple form.

\begin{definition}
    A formula in a relational language is called \emph{simple} if it is quantifier-free. Equivalently, it is a finite Boolean combination of atomic formulas.

    Let $\bfa=\langle \bfa_0,\prec^A\rangle$ be an $L$-structure and let $\varphi(x,y)$ be a simple formula in the language $L$. We write
    \[
        \varphi^{\bfa}(a,b)
    \]
    to mean that $\bfa\models \varphi(a,b)$.
\end{definition}

\begin{definition}
    \label{def:simple-bi-definability}
    Let $\k$ be a class of finite $L$-structures and let $\k'$ be a class of finite $L'$-structures. We say that $\k$ and $\k'$ are \emph{simply bi-definable} if there exist simple formulas
    \[
        \varphi(x,y)\ \text{ in } L
        \qquad\text{and}\qquad
        \varphi'(x,y)\ \text{ in } L'
    \]
    such that the following hold.

    For every
    \[
        \bfa=\langle \bfa_0,\prec^A\rangle\in \k,
    \]
    if we define a binary relation $<'^{\,A}$ on $A$ by
    \[
        a<'^{\,A} b
        \quad\Longleftrightarrow\quad
        \bfa\models \varphi(a,b),
    \]
    then
    \[
        \langle \bfa_0,<'^{\,A}\rangle\in \k'.
    \]

    Similarly, for every
    \[
        \bfb=\langle \bfb_0,\prec'^{\,B}\rangle\in \k',
    \]
    if we define a binary relation $<^{B}$ on $B$ by
    \[
        a<^{B} b
        \quad\Longleftrightarrow\quad
        \bfb\models \varphi'(a,b),
    \]
    then
    \[
        \langle \bfb_0,<^{B}\rangle\in \k.
    \]

    Moreover, these two constructions are inverse to one another up to isomorphism.
\end{definition}

In other words, $\k$ and $\k'$ are simply bi-definable if the two orderings can be recovered from one another by quantifier-free formulas over the common reduct language.

The relevance of this notion is that it preserves the two combinatorial properties that were central in Chapters~5 and~8.

\begin{proposition}
    \label{prop:simple-bi-definability-preserves-properties}
    Let $\k$ and $\k'$ be simply bi-definable classes of finite ordered structures with the same reduct class $\k_0$. Then:
    \begin{enumerate}
        \item[(i)] $\k$ has the Ramsey property if and only if $\k'$ has the Ramsey property;
        \item[(ii)] $\k$ has the ordering property if and only if $\k'$ has the ordering property.
    \end{enumerate}
\end{proposition}

\begin{proof}
    We prove the forward implications; the converse directions are symmetric.

    Assume first that $\k$ has the Ramsey property. Let
    \[
        \bfa'\leq \bfb'
    \]
    be structures in $\k'$, and let $k\geq 2$. Using simple bi-definability, let $\bfa,\bfb\in\k$ be the corresponding structures on the same underlying sets. Choose
    \[
        \bfc\in\k
    \]
    such that
    \[
        \bfc\to (\bfb)^{\bfa}_k.
    \]
    Let $\bfc'\in\k'$ be the structure corresponding to $\bfc$ under the defining formula $\varphi$. Since embeddings in one class are exactly embeddings in the other after translating the order relation by $\varphi$ and $\varphi'$, the partition relation for $\bfc$ transfers to the corresponding partition relation for $\bfc'$. Hence
    \[
        \bfc'\to (\bfb')^{\bfa'}_k.
    \]
    Thus $\k'$ has the Ramsey property.

    Now assume that $\k$ has the ordering property. Let
    \[
        \bfa_0\in\k_0.
    \]
    Choose an expansion
    \[
        \bfa'=\langle \bfa_0,\prec'\rangle\in\k',
    \]
    and let $\bfa\in\k$ be the corresponding structure. By the ordering property for $\k$, there exists
    \[
        \bfb_0\in\k_0
    \]
    such that every member of $\k$ with reduct $\bfb_0$ contains a copy of $\bfa$. Let
    \[
        \bfb'=\langle \bfb_0,\triangleleft\rangle\in\k'
    \]
    be arbitrary, and let $\bfb\in\k$ be the corresponding structure. Then
    \[
        \bfa\leq \bfb,
    \]
    and therefore, translating back by simple bi-definability,
    \[
        \bfa'\leq \bfb'.
    \]
    Thus $\k'$ has the ordering property.
\end{proof}

We now come to the main uniqueness theorem.

\begin{theorem}
    \label{thm:uniqueness-simple-bi-definability}
    Let $\k$ and $\k'$ be reasonable Fra\"iss\'e order classes in the signatures $L$ and $L'$, and suppose that they have the same reduct
    \[
        \k_0.
    \]
    Assume that both $\k$ and $\k'$ have the Ramsey property and the ordering property. Then $\k$ and $\k'$ are simply bi-definable.
\end{theorem}

\begin{proof}
    Let
    \[
        \bff=\langle \bff_0,\prec_0\rangle=\mathrm{Flim}(\k)
        \qquad\text{and}\qquad
        \bff'=\langle \bff_0,\prec_0'\rangle=\mathrm{Flim}(\k'),
    \]
    where the two limits are taken on the same reduct $\bff_0=\mathrm{Flim}(\k_0)$, using Chapter~7.

    Let
    \[
        G_0=\Aut(\bff_0),
        \qquad
        X_\k=\overline{G_0\cdot \prec_0},
        \qquad
        X_{\k'}=\overline{G_0\cdot \prec_0'}.
    \]
    By Chapter~8, since both $\k$ and $\k'$ have the Ramsey property and the ordering property, the flows $X_\k$ and $X_{\k'}$ are both universal minimal flows of $G_0$. Hence there exists a $G_0$-flow isomorphism
    \[
        \pi\colon X_\k\to X_{\k'}.
    \]

    We claim that
    \[
        \pi(G_0\cdot \prec_0)=G_0\cdot \prec_0'.
    \]
    Indeed, for a linear ordering $\prec\in X_\k$, the condition
    \[
        \prec\in G_0\cdot \prec_0
    \]
    is equivalent to the statement that the ordered structure
    \[
        \langle \bff_0,\prec\rangle
    \]
    is isomorphic to $\bff$, or equivalently that it has age $\k$ and the extension property. This witnesses the triangle condition for $\k'$. Applying Corollary~\ref{cor:triangle-condition} with the roles of $\k$ and $\k'$ interchanged, we obtain
    \[
        \k=\k'
        \qquad\text{or}\qquad
        \k=(\k')^*.
    \]
    \begin{mynote}
        Something is missing. what the heck is \(\k^*\)?
    \end{mynote}
    Equivalently,
    \[
        \k'=\k
        \qquad\text{or}\qquad
        \k'=\k^*.
    \]
    Similarly, $G_0\cdot \prec_0'$ is a dense $G_\delta$ subset of $X_{\k'}$. Because $\pi$ is a homeomorphism, it must send the unique dense $G_\delta$ orbit of $X_\k$ onto the unique dense $G_\delta$ orbit of $X_{\k'}$. This proves the claim.

    Composing $\pi$ with a suitable element of $G_0$ if necessary, we may therefore assume that
    \[
        \pi(\prec_0)=\prec_0'.
    \]
    Let
    \[
        G=\Aut(\bff)
        \qquad\text{and}\qquad
        G'=\Aut(\bff').
    \]
    Then
    \[
        G=\Stab_{G_0}(\prec_0)
        \qquad\text{and}\qquad
        G'=\Stab_{G_0}(\prec_0').
    \]
    Since $\pi$ is $G_0$-equivariant and sends $\prec_0$ to $\prec_0'$, it follows that
    \[
        G=G'.
    \]

    We now define $<'$ from $<$ by a simple formula. Since $G=G'$, the relation
    \[
        x<' y
    \]
    is invariant under $\Aut(\bff)$. Because $\bff$ is ultrahomogeneous in the relational language $L$, every $\Aut(\bff)$-invariant binary relation is a union of $2$-orbits, and each such orbit is determined by the quantifier-free type of the ordered pair in $\bff$. There are only finitely many such $2$-orbits, so the relation $<'$ is defined in $\bff$ by a simple formula
    \[
        \varphi(x,y)
    \]
    in the language $L$.

    By the same argument, the relation $<$ is defined in $\bff'$ by a simple formula
    \[
        \varphi'(x,y)
    \]
    in the language $L'$.

    Finally, let
    \[
        \bfa=\langle \bfa_0,\prec^A\rangle\in\k.
    \]
    Since $\k=\Age(\bff)$, the structure $\bfa$ embeds into $\bff$. Restricting the formula $\varphi$ to the image of such an embedding defines on $A$ a second linear order $<'^{\,A}$, and the resulting structure
    \[
        \langle \bfa_0,<'^{\,A}\rangle
    \]
    belongs to $\Age(\bff')=\k'$. The same reasoning applies in the other direction using $\varphi'$. Because both formulas come from the same subgroup
    \[
        G=G'\leq G_0,
    \]
    the two constructions are inverse to one another up to isomorphism. Therefore $\k$ and $\k'$ are simply bi-definable.
\end{proof}

\begin{remark}
    If the signatures are finite, the same argument yields first-order bi-definability in the usual sense. For the purposes of the present thesis, however, the quantifier-free version is the natural one and is sufficient.
\end{remark}

The theorem shows that a reasonable Ramsey expansion with the ordering property is unique up to a very rigid notion of definitional equivalence. In the next section we record some strengthenings and examples of this phenomenon.

\section{Strengthenings and examples}

We now record two useful strengthenings of Theorem~\ref{thm:uniqueness-simple-bi-definability}, together with some examples.

\begin{proposition}
    \label{prop:comparable-expansions-equal}
    In the context of Theorem~\ref{thm:uniqueness-simple-bi-definability}, suppose that $\k$ and $\k'$ are comparable under inclusion, that is,
    \[
        \k\subseteq \k'
        \qquad\text{or}\qquad
        \k'\subseteq \k.
    \]
    Then
    \[
        \k=\k'.
    \]
    In particular, if
    \[
        \k=\k_0*\mathcal{LO},
    \]
    then $\k=\k'$.
\end{proposition}

\begin{proof}
    Assume, for example, that
    \[
        \k\subseteq \k'.
    \]
    Suppose toward a contradiction that $\k\neq\k'$. Then there exists
    \[
        \bfa=\langle \bfa_0,\prec\rangle\in \k'\setminus \k.
    \]
    Since $\k'$ has the ordering property, there exists
    \[
        \bfb_0\in \k_0
    \]
    such that for every expansion
    \[
        \bfb'=\langle \bfb_0,\prec'\rangle\in \k'
    \]
    one has
    \[
        \bfa\leq \bfb'.
    \]
    Choose any expansion
    \[
        \bfb=\langle \bfb_0,\prec'\rangle\in \k.
    \]
    Since $\k\subseteq \k'$, we also have $\bfb\in\k'$, and therefore
    \[
        \bfa\leq \bfb.
    \]
    But $\k$ is hereditary, so this implies
    \[
        \bfa\in \k,
    \]
    a contradiction. Hence $\k=\k'$.

    For the final statement, note that if
    \[
        \k=\k_0*\mathcal{LO},
    \]
    then every expansion in $\k'$ also belongs to $\k$, so
    \[
        \k'\subseteq \k.
    \]
    The first part of the proposition therefore gives $\k=\k'$.
\end{proof}

Thus in all the free-expansion examples from Chapter~9, the expansion with both the Ramsey and ordering properties is uniquely determined.

There is another natural way in which two expansions may differ. If \(\k\) is a class of ordered structures, let
\[
    \k^*=\{\langle \bfa_0,\prec^*\rangle : \langle \bfa_0,\prec\rangle\in \k\},
\]
where
\[
    a\prec^* b
    \quad\Longleftrightarrow\quad
    b\prec a.
\]
Then $\k^*$ is simply bi-definable with $\k$, so by Proposition~\ref{prop:simple-bi-definability-preserves-properties}, the classes $\k$ and $\k^*$ either both have, or both fail to have, the Ramsey and ordering properties. The next condition shows that in some cases these are the only possibilities.

Given
\[
    \bfa=\langle \bfa_0,\prec\rangle\in\k
    \qquad\text{and}\qquad
    a,b\in A_0
\]
with
\[
    a\prec b,
\]
we write
\[
    \operatorname{tp}_{\bfa}(a,b)
\]
for the atomic type of the ordered pair $(a,b)$ in the language $L$. Equivalently, this is the quantifier-free type of $(a,b)$ in $\bfa$.

\begin{definition}
    \label{def:triangle-condition}
    We say that $\k$ satisfies the \emph{triangle condition} if for any two distinct pair-types
    \[
        \sigma\neq \tau
    \]
    realized in members of $\k$, there exist
    \[
        \bfa=\langle \bfa_0,\prec\rangle\in \k
        \qquad\text{and}\qquad
        a\prec b\prec c
    \]
    in $A_0$ such that either
    \[
        \operatorname{tp}_{\bfa}(a,b)=\operatorname{tp}_{\bfa}(b,c)=\sigma
        \qquad\text{and}\qquad
        \operatorname{tp}_{\bfa}(a,c)=\tau,
    \]
    or
    \[
        \operatorname{tp}_{\bfa}(a,b)=\operatorname{tp}_{\bfa}(b,c)=\tau
        \qquad\text{and}\qquad
        \operatorname{tp}_{\bfa}(a,c)=\sigma.
    \]
\end{definition}

The point is that the triangle condition forces any alternative admissible ordering to agree either with the original ordering or with its reverse.

\begin{corollary}
    \label{cor:triangle-condition}
    In the context of Theorem~\ref{thm:uniqueness-simple-bi-definability}, assume that $\k$ satisfies the triangle condition. Then
    \[
        \k'=\k
        \qquad\text{or}\qquad
        \k'=\k^*.
    \]
\end{corollary}

\begin{proof}
    By Theorem~\ref{thm:uniqueness-simple-bi-definability}, the classes $\k$ and $\k'$ are simply bi-definable. In the notation of its proof, we may therefore realize the two Fra\"iss\'e limits on the same reduct
    \[
        \bff_0
    \]
    and write
    \[
        \bff=\langle \bff_0,\prec_0\rangle,
        \qquad
        \bff'=\langle \bff_0,\prec_0'\rangle,
    \]
    with
    \[
        \Aut(\bff)=\Aut(\bff').
    \]
    It is enough to show that
    \[
        \prec_0'=\prec_0
        \qquad\text{or}\qquad
        \prec_0'=\prec_0^*.
    \]

    Assume this fails. Then there exist
    \[
        a,b,c,d\in F_0
    \]
    such that
    \[
        a\prec_0 b,\qquad c\prec_0 d,
    \]
    but
    \[
        a\prec_0' b,\qquad d\prec_0' c.
    \]
    Let
    \[
        \sigma=\operatorname{tp}_{\bff}(a,b)
        \qquad\text{and}\qquad
        \tau=\operatorname{tp}_{\bff}(c,d).
    \]
    Then
    \[
        \sigma\neq \tau,
    \]
    since $\prec_0'$ is invariant under $\Aut(\bff)=\Aut(\bff')$.

    By the triangle condition, there exist
    \[
        \bfa\in\k
        \qquad\text{and}\qquad
        u\prec v\prec w
    \]
    in $A_0$ such that either
    \[
        \operatorname{tp}_{\bfa}(u,v)=\operatorname{tp}_{\bfa}(v,w)=\sigma
        \quad\text{and}\quad
        \operatorname{tp}_{\bfa}(u,w)=\tau,
    \]
    or the same statement with $\sigma$ and $\tau$ interchanged. Realizing $\bfa$ as a substructure of $\bff$, the relation $\prec_0'$ must orient the pairs of type $\sigma$ as in $(a,b)$ and the pairs of type $\tau$ as in $(d,c)$. In the first case this gives
    \[
        u\prec_0' v\prec_0' w
        \qquad\text{and}\qquad
        w\prec_0' u,
    \]
    a contradiction; in the second case one obtains
    \[
        w\prec_0' v\prec_0' u
        \qquad\text{and}\qquad
        u\prec_0' w,
    \]
    again a contradiction. Therefore
    \[
        \prec_0'=\prec_0
        \qquad\text{or}\qquad
        \prec_0'=\prec_0^*,
    \]
    and hence
    \[
        \k'=\k
        \qquad\text{or}\qquad
        \k'=\k^*.
    \]
\end{proof}

We now record two examples.

\begin{example}
    Let $\k_0$ be the class of all finite equivalence relations, and let $\k$ be the class of all finite convexly ordered equivalence relations. Then $\k$ is the unique reasonable Fra\"iss\'e order expansion of $\k_0$ with both the Ramsey property and the ordering property.
\end{example}

\begin{proof}
    Let $\k'$ be any reasonable Fra\"iss\'e order expansion of $\k_0$ with both properties. If
    \[
        \k'\subseteq \k,
    \]
    then Proposition~\ref{prop:comparable-expansions-equal} yields
    \[
        \k'=\k.
    \]
    So suppose instead that $\k'\nsubseteq \k$. Then $\k'$ contains some
    \[
        \bfa'=\langle \bfa_0,E,\prec'\rangle
    \]
    with
    \[
        a\prec' b\prec' c,
        \qquad
        a\,E\,c,
        \qquad
        a\not E b,
        \qquad
        b\not E c,
    \]
    for some $a,b,c\in A_0$. This witnesses the triangle condition for $\k'$, and therefore Corollary~\ref{cor:triangle-condition} implies that
    \[
        \k'=\k
        \qquad\text{or}\qquad
        \k'=\k^*.
    \]
    But $\k^*=\k$ in this case, since reversing a convex ordering is again convex. Hence
    \[
        \k'=\k.
    \]
\end{proof}

\begin{example}
    Let $\k_0$ be the class of all finite posets, and let $\k$ be the class of all finite structures
    \[
        \langle A,\sqsubset,\prec\rangle
    \]
    where $\langle A,\sqsubset\rangle$ is a poset and $\prec$ is a linear ordering extending $\sqsubset$. Then the only reasonable Fra\"iss\'e order expansions of $\k_0$ with both the Ramsey and ordering properties are $\k$ and $\k^*$.
\end{example}

\begin{proof}
    This class has the Ramsey and ordering properties. It also satisfies the triangle condition: one may realize two distinct pair-types by appropriately choosing a three-element chain
    \[
        a\prec b\prec c
    \]
    and using the compatibility of $\prec$ with the partial order. Corollary~\ref{cor:triangle-condition} therefore gives the conclusion.
\end{proof}

The situation is not always so rigid. In some classes there are several non-isomorphic expansions satisfying the Ramsey and ordering properties. This occurs, for example, for finite Boolean algebras and for finite-dimensional vector spaces over a finite field. Since the classification of all such expansions would take us too far from the main line of the thesis, we do not pursue it here.

The next section explains the conceptual meaning of the uniqueness theorem from the viewpoint of minimality among Ramsey expansions.

\section{Minimality among Ramsey expansions}

We now explain the sense in which expansions with both the Ramsey property and the ordering property are minimal among Ramsey expansions of a fixed reduct class.

\begin{theorem}
    \label{thm:minimality-among-ramsey-expansions}
    Let $\k_0$ be a Fra\"iss\'e class in a relational signature $L_0$, let
    \[
        L=L_0\cup\{<\},
    \]
    and let $\k,\k'$ be reasonable Fra\"iss\'e order classes in the language $L$ expanding $\k_0$. Assume that $\k$ has the Ramsey property and that $\k'$ has both the Ramsey property and the ordering property. Then there exists a reasonable Fra\"iss\'e order class $\k''$ in the language $L$ such that
    \[
        \k''\subseteq \k
    \]
    and $\k''$ is simply bi-definable with $\k'$.
\end{theorem}

\begin{proof}
    Let
    \[
        \bff=\langle \bff_0,\prec_0\rangle=\mathrm{Flim}(\k),
        \qquad
        \bff'=\langle \bff_0,\prec_0'\rangle=\mathrm{Flim}(\k'),
    \]
    where we identify the two reduct limits using Chapter~7. Put
    \[
        G_0=\Aut(\bff_0),
        \qquad
        G=\Aut(\bff).
    \]
    Let
    \[
        X_\k=\overline{G_0\cdot \prec_0}
        \qquad\text{and}\qquad
        X_{\k'}=\overline{G_0\cdot \prec_0'}.
    \]

    Since $\k$ has the Ramsey property, Corollary~\ref{cor:fraisse-order-class-ea} shows that
    \[
        G=\Aut(\bff)
    \]
    is extremely amenable. Let $Y$ be any minimal subflow of $X_\k$. We claim that $Y$ is the universal minimal flow of $G_0$.

    Indeed, let $Z$ be any minimal compact $G_0$-flow. Restricting the action to the subgroup $G$, extreme amenability yields a point
    \[
        z_0\in Z
    \]
    fixed by every element of $G$. By Theorem~\ref{thm:ramsey-universal-ambit}, there exists a surjective $G_0$-map
    \[
        \pi\colon X_\k\to Z.
    \]
    Since $Y$ is a minimal subflow of $X_\k$, the image $\pi(Y)$ is a nonempty closed $G_0$-invariant subset of $Z$, hence
    \[
        \pi(Y)=Z.
    \]
    Thus every minimal compact $G_0$-flow is a factor of $Y$, so $Y$ is universal minimal.

    On the other hand, since $\k'$ has both the Ramsey property and the ordering property, Corollary~\ref{cor:ramsey-ordering-umf} implies that
    \[
        X_{\k'}
    \]
    is also the universal minimal flow of $G_0$. Therefore there is a $G_0$-flow isomorphism
    \[
        \theta\colon X_{\k'}\to Y.
    \]
    Replacing $\theta$ by $g\cdot \theta$ for a suitable $g\in G_0$ if necessary, we may assume that
    \[
        \theta(\prec_0')=\prec_0''
        \in Y.
    \]

    Since $\theta$ is $G_0$-equivariant, the stabilizers of $\prec_0'$ and $\prec_0''$ in $G_0$ coincide. Hence
    \[
        \Aut(\langle \bff_0,\prec_0'\rangle)
        =
        \Aut(\langle \bff_0,\prec_0''\rangle).
    \]
    Exactly as in the proof of Theorem~\ref{thm:uniqueness-simple-bi-definability}, it follows that there are simple formulas defining $\prec_0'$ from $\prec_0''$ and $\prec_0''$ from $\prec_0'$. In particular,
    \[
        \langle \bff_0,\prec_0''\rangle
    \]
    is ultrahomogeneous.

    Let
    \[
        \k''=\Age(\langle \bff_0,\prec_0''\rangle).
    \]
    Then $\k''$ is a reasonable Fra\"iss\'e order class expanding $\k_0$, and by construction it is simply bi-definable with $\k'$.

    It remains to show that
    \[
        \k''\subseteq \k.
    \]
    Since
    \[
        \prec_0''\in Y\subseteq X_\k,
    \]
    Proposition~\ref{prop:admissible-orderings-characterization} shows that every finite substructure of
    \[
        \langle \bff_0,\prec_0''\rangle
    \]
    belongs to $\k$. Hence
    \[
        \Age(\langle \bff_0,\prec_0''\rangle)\subseteq \k,
    \]
    that is,
    \[
        \k''\subseteq \k.
    \]
    This completes the proof.
\end{proof}

Thus, up to simple bi-definability, the expansions of a given reduct class that have both the Ramsey property and the ordering property are the smallest among those that have the Ramsey property. The proof also suggests a natural strategy: starting from a Ramsey expansion $\k$, one may try to pass to a minimal subflow of $X_\k$ and recover from it a smaller expansion with better dynamical properties.

In the next chapter this idea is developed further through Ramsey degrees, which allow one to treat situations in which no reasonable expansion with the full Ramsey property is available.